\documentclass[12pt]{article}

%% Packages
\usepackage[margin=1.25in]{geometry}
\usepackage{setspace}
\usepackage{hyperref}
\usepackage{natbib}
\usepackage{amsmath}
\usepackage{microtype}
\usepackage{booktabs}
\usepackage{xcolor}

\hypersetup{
  colorlinks=true,
  linkcolor=blue!60!black,
  citecolor=blue!60!black,
  urlcolor=blue!60!black
}

\onehalfspacing

\title{\textbf{From Bond-Backed Currency to Real Bills:\\
The Inelasticity Problem and the Origins of the Federal Reserve}}

\author{Thomas J.\ Sargent}
\date{March 2026}

\begin{document}

\maketitle

\begin{abstract}
\noindent
The National Banking Acts of 1863--1864 created federally chartered banks whose
circulating notes had to be backed by holdings of specified U.S.\ government bonds
deposited with the Comptroller of the Currency.  Critics from the 1870s onward
complained that this arrangement made the currency ``inelastic''---unresponsive to
the seasonal and cyclical needs of trade.  Over the following four decades a range
of reform proposals sought to broaden the permissible backing to include railroad
bonds, agricultural warehouse receipts, and other non-government assets.  The
Aldrich--Vreeland Act of 1908, passed in the wake of the Panic of~1907, temporarily
allowed emergency notes backed by a wider range of securities and created the
National Monetary Commission that would lay the intellectual groundwork for the
Federal Reserve.  When the Federal Reserve Act was finally passed in December~1913,
its long title explicitly promised ``to furnish an elastic currency, to afford means
of rediscounting commercial paper.''  The Act fulfilled this promise by embedding
the \emph{real bills doctrine} into the Fed's discount window: only short-term,
self-liquidating commercial paper---so-called ``real bills''---would be eligible
for rediscount.  This essay traces the thread from the Civil War banking statutes
through the inelasticity debate, the Aldrich--Vreeland interlude, and finally to
the real bills provisions of the 1913 Act, together with the subsequent macroeconomic
controversies those provisions engendered.
\end{abstract}

\tableofcontents
\newpage

%% ------------------------------------------------------------------
\section{The National Banking System, 1863--1864}
\label{sec:national-banking}
%% ------------------------------------------------------------------

\subsection{Origins in the Civil War}

Before the Civil War the United States had no uniform national currency and no
federal bank-chartering authority.  After the demise of the Second Bank of the
United States in 1836, banking regulation had devolved almost entirely to the states.
The resulting system was fragmented: notes issued by one state's banks circulated at a
steep discount elsewhere, and ``wildcat banking'' scandals in states with lax rules
had badly shaken public confidence in paper money.

The exigencies of financing the Union war effort changed the political calculus
decisively.  Secretary of the Treasury Salmon P.\ Chase, who had opposed federal
involvement in banking before the war, came to see a national currency as
indispensable both for fiscal and monetary reasons.  Congress responded with the
\emph{National Currency Act} of February~25, 1863 (12~Stat.\ 665, Ch.\ 58), later
renamed the \emph{National Bank Act}, and its more comprehensive successor of
June~3, 1864 (13~Stat.\ 99, Ch.\ 106).

\subsection{The Bond-Backing Requirement}

The central mechanism of the new system was straightforward.  A bank wishing to
issue National Bank Notes had to:
\begin{enumerate}
  \item Obtain a federal charter from the newly created \emph{Office of the
    Comptroller of the Currency} (OCC), housed in the Treasury Department.
  \item Deposit with the Comptroller a specified quantity of United States
    government bonds---the 1863~Act required that no less than one-third of the
    bank's capital take the form of such bonds, held in the Comptroller's vault.
  \item Receive in return printed currency (National Bank Notes) up to 90~percent
    of the par or market value (whichever was lower) of the bonds so deposited.
\end{enumerate}

The notes themselves were uniform in design, printed by the federal government, and
identical in appearance across all issuing banks save for a small identification
panel.  The Comptroller was charged with inspecting banks to verify compliance with
the holding requirement---making the OCC the first systematic federal bank
supervisor in American history.  Hugh McCulloch, the first Comptroller
(1863--1865), personally evaluated charter applications, oversaw note design, and
arranged for note engraving, printing, and distribution.

State bank notes came under federal attack.  A 10~percent federal tax
imposed in 1865--1866 on payments made in state bank notes effectively drove them
out of circulation, leaving National Bank Notes (and residual Civil War
greenbacks) as the dominant hand-to-hand currency.  The number of national banks
rose from 66 immediately after the 1863~Act to 7{,}473 by 1913; the number of state
banks fell from 1{,}466 in 1863 to 247 by 1868, before rebounding as bankers
exploited the lower capital requirements of state charters and the growing
popularity of demand deposits (which were not taxed).

\subsection{The Comptroller as Enforcer}

That the U.S.\ Comptroller of the Currency was established in large part to inspect
and verify compliance with the bond-deposit requirement deserves emphasis.  The
Comptroller examined national banks to confirm that the bonds pledged as backing for
outstanding notes were actually present and of sufficient value.  This supervisory
function was novel: previously there was no federal authority empowered to walk
into a bank and verify its balance sheet.  For our purposes the key point is that
the legitimacy of every dollar of National Bank Note currency depended on
there being a corresponding U.S.\ government bond sitting in the Comptroller's
vault in Washington.

%% ------------------------------------------------------------------
\section{The Inelasticity Problem, 1865--1907}
\label{sec:inelasticity}
%% ------------------------------------------------------------------

\subsection{Why Bond-Backing Created Rigidity}

The very feature that gave National Bank Notes their credibility---the requirement
that each dollar of notes be backed by a federal bond---also made the currency
supply hostage to fiscal policy rather than to the needs of commerce.  The total
stock of National Bank Notes that could lawfully circulate was bounded above by
(roughly) 90~percent times the total par value of the designated government bonds held
by all national banks in aggregate.  Three structural features made this constraint
binding:

\begin{enumerate}
  \item \textbf{Declining bond supply after the Civil War.}  The federal government
    ran large primary surpluses from the mid-1860s through the 1880s and used the
    proceeds to retire its war debt.  As the supply of outstanding eligible bonds
    fell, the monetary base of National Bank Notes contracted or was capped even as
    the economy grew.

  \item \textbf{Premium prices for eligible bonds.}  Because holding eligible bonds
    entitled a bank to issue profitable currency, demand for those bonds pushed
    their prices \emph{above} par.  A bank that paid a premium for a bond received
    notes at the bond's par value, not its market price, compressing the profit
    margin of note issue and reducing the incentive to expand note circulation.

  \item \textbf{Seasonal and geographic rigidities.}  The agrarian economy of the
    South and West needed sharply higher volumes of currency during planting and
    harvest seasons.  Because National Bank Note issuance was tied to bond
    portfolios---not to seasonal credit demand---banks could not expand circulation
    when farmers needed cash and contract it again when crops were sold.  The result
    was recurrent seasonal currency shortages and elevated interest rates in interior
    agricultural districts, punctuated by periodic financial panics.
\end{enumerate}

Friedman and Schwartz \citep{friedman-schwartz1963} documented that the currency
component of the money supply under the National Banking System was indeed
far less responsive to economic conditions than one would expect from a well-designed
monetary system.  Johnson \citep{johnson1999} summarized the pre-Fed banking system
as characterized by ``immobile reserves and inelastic currency.''

\subsection{Political Demands for Currency Reform}

The inelasticity of the currency became a major political grievance from the
1870s onward.  Broadly, two coalitions formed:

\paragraph{Greenbackers and silver inflationists.}  The most radical critics demanded
that the federal government simply issue more fiat money---greenbacks---or that
silver be monetized at a generous ratio.  William Jennings Bryan's famous ``Cross
of Gold'' speech at the 1896 Democratic Convention crystallized this position.
These advocates cared less about the \emph{backing} of currency than about its
\emph{quantity}; they wanted the money supply to be large enough to relieve
debt-burdened farmers and spur economic activity.

\paragraph{Asset-currency reformers.}  A more technically focused group, drawn
largely from commercial banking circles in the Northeast but with agrarian
sympathizers in the Midwest, argued that the problem was specifically the
\emph{bond-backed} structure of note issue.  Their proposed remedy was
\emph{asset currency}: allowing national banks to issue notes backed not only by
government bonds but also by a wider range of productive assets.  The menu of
proposed alternative backing included:
\begin{itemize}
  \item \emph{Railroad and corporate bonds.}  By the 1880s railroad bonds vastly
    exceeded federal bonds in total outstanding volume; asset-currency advocates
    argued that high-grade railroad paper was just as safe and far more plentiful.
  \item \emph{Agricultural warehouse receipts.}  Farmers' groups and some midwestern
    bankers proposed that banks should be able to issue notes against warehouse
    receipts for stored grain, cotton, and other commodities.  This would make
    the currency supply directly responsive to the seasonal rhythm of agriculture.
  \item \emph{General ``good assets.''} Some proposals envisaged allowing note
    issue against any sound asset on a bank's books, subject to regulatory
    approval---a forerunner of the broader ``eligible paper'' concept that would
    enter the Federal Reserve Act.
\end{itemize}

The common thread in asset-currency proposals was the complaint that the existing
system was ``too inelastic,'' that it prevented currency from being ``responsive to
the needs of trade.''  This phrase---\emph{responsiveness to the needs of trade}---
would recur across decades of monetary debate and find its ultimate echo in both
the Aldrich--Vreeland Act and the Federal Reserve Act.

\subsection{Notable Reform Proposals and Legislation}

Successive Secretaries of the Treasury and various congressional committees floated
asset-currency bills through the 1880s and 1890s.  The Indianapolis Monetary
Convention of 1897 brought together business and banking representatives who formally
endorsed asset currency backed by commercial paper.  The American Bankers
Association lobbied actively for a reform along these lines throughout the 1890s
and early 1900s.  None of these proposals passed Congress, stymied by the
cross-cutting politics of money: Eastern bankers feared that broadening the
acceptable backing would destabilize the currency, while agrarians found
asset-currency reform insufficiently generous.

%% ------------------------------------------------------------------
\section{The Panic of 1907 and the Aldrich--Vreeland Act of 1908}
\label{sec:aldrich-vreeland}
%% ------------------------------------------------------------------

\subsection{The Panic of 1907}

The Panic of 1907 brought the defects of the National Banking System into sharp
relief.  A stock-market crash and a run on several New York trust companies
triggered a cascade of bank suspensions.  The crisis was arrested only through the
personal intervention of J.P.\ Morgan, who organized a pool of private funds to
shore up the threatened institutions---a striking demonstration that in the
absence of a lender of last resort, private coordination by a single wealthy
financier was the banking system's only circuit-breaker.  The episode produced a broad
consensus across political lines that the banking and currency system needed
fundamental reform.

\subsection{The Aldrich--Vreeland Act (May 30, 1908)}

Congress responded with unusual speed.  The Aldrich--Vreeland Act (Pub.\ L.\
60-169, 35~Stat.\ 546), introduced by Representative Edward Vreeland
(Republican, New York) and championed in the Senate by Senator Nelson W.\ Aldrich
(Republican, Rhode Island), was signed into law by President Theodore Roosevelt on
May 30, 1908.

The Act had two main provisions.

\paragraph{Emergency currency.}  National banks could form \emph{National Currency
Associations} in groups of ten or more, with aggregate capital of at least
\$5~million.  Having formed such an association, member banks could issue
\emph{emergency currency} backed not merely by government bonds but by
``almost any securities the banks were holding''---including railroad and other
corporate bonds.  Once approved by association officers and the Comptroller of
the Currency, these notes could be placed in circulation as legal tender.  The
emergency notes bore a tax of 5~percent per annum for the first month outstanding,
rising by 1~percentage point each subsequent month.  This punitive tax was designed
to ensure that the emergency currency would be retired promptly once conditions
normalized.  As it turned out, no notes were issued under this provision during
peacetime conditions.

\paragraph{National Monetary Commission.}  The Act created the \emph{National
Monetary Commission} (NMC), chaired by Senator Aldrich, to study banking and
currency systems in the United States and abroad and to recommend comprehensive
reforms.  The NMC commissioned several dozen studies of banking practices in
Europe, particularly in France, Germany, and Britain, and its reports became the
intellectual foundation on which the Federal Reserve System was constructed.

\subsection{The Aldrich--Vreeland Act in Practice: August--October 1914}

The emergency currency provisions proved their worth at the outbreak of World War~I
in August~1914, when foreign creditors---particularly British banks---demanded
immediate gold payments that would otherwise have drained the U.S.\ monetary base.
Secretary of the Treasury William Gibbs McAdoo activated the Aldrich--Vreeland
machinery, and by October~23, 1914, some \$368{,}616{,}990 in emergency notes had
been placed in circulation, averting a banking panic.  The notes were entirely
withdrawn once the Federal Reserve banks opened for business in November~1914.

\subsection{From Aldrich--Vreeland to the Federal Reserve}

The NMC submitted its final report to Congress on January~9, 1912.  The report
accompanied draft legislation---the \emph{Aldrich Plan}---proposing a
centralized National Reserve Association with fifteen regional branches.  The Plan
called for a lender of last resort that would discount commercial paper for member
banks, making the currency supply elastic and eliminating the bond-backing
straitjacket.  Opposition arose quickly: rural and progressive Democrats feared
that the Association would be controlled by Wall Street's ``Money Trust.''  After
extensive hearings by the Pujo Committee (May~1912--January~1913) into alleged
financial concentration, President-elect Wilson's administration stepped in with
a modified plan that preserved the decentralization demanded by the Bryan wing of
the Democratic Party.

%% ------------------------------------------------------------------
\section{The Federal Reserve Act of 1913 and the Real Bills Doctrine}
\label{sec:federal-reserve}
%% ------------------------------------------------------------------

\subsection{The Long Title as a Statement of Intent}

The Federal Reserve Act, signed by President Woodrow Wilson on December~23, 1913,
bore the official long title:
\begin{quote}
\emph{An Act to provide for the establishment of Federal reserve banks, to furnish
an elastic currency, to afford means of rediscounting commercial paper, to
establish a more effective supervision of banking in the United States, and for
other purposes.}
\end{quote}

The phrase ``to furnish an elastic currency'' was a direct answer to the
four-decade complaint about the National Banking System.  ``To afford means of
rediscounting commercial paper'' signaled the mechanism by which elasticity would be
achieved: member banks could bring eligible commercial paper to their regional
Federal Reserve bank, which would discount it (lend against it), providing the
bank with reserves and thus allowing it to expand its loans and note issue.

\subsection{Eligible Paper and the Real Bills Doctrine}

What made paper ``eligible'' for rediscount at the Fed?  Section~13 of the
Act specified that the Federal Reserve banks could discount ``notes, drafts, and
bills of exchange arising out of actual commercial transactions.''  The key
adjective was \emph{commercial}: the paper had to represent a real transaction in
goods or services, not a speculative financial transaction.  Further criteria
required that eligible paper be:
\begin{itemize}
  \item \emph{Short-term}: maturing within 90 days (or six months for agricultural
    paper);
  \item \emph{Self-liquidating}: arising from a specific productive transaction
    whose completion would generate the funds to repay the loan; and
  \item \emph{Drawn for a commercial purpose}: explicitly excluding paper ``issued
    or drawn for the purpose of carrying or trading in stocks, bonds, or other
    investment securities.''
\end{itemize}

This set of criteria precisely codified what monetary theorists called the
\emph{real bills doctrine} (also known as the \emph{commercial loan theory of
banking}).

\subsection{Intellectual Origins of the Real Bills Doctrine}

The real bills doctrine has roots in seventeenth- and eighteenth-century banking
thought.  John Law (\citeyear{law1720}) sketched the core idea---an
``output-governed currency secured to real property and responding to the needs of
trade.''  Adam Smith gave the doctrine its classical formulation in \emph{The
Wealth of Nations} (Book~II, Ch.~2), arguing that a bank that confines its loans to
the discount of real bills drawn against actual goods in production will never
over-issue its notes: as goods are sold and bills repaid, the notes return to the
bank and are extinguished, so the money supply automatically tracks the real volume
of transactions.

Henry Thornton (\citeyear{thornton1802}) identified the doctrine's central flaw
in his \emph{An Enquiry into the Nature and Effects of the Paper Credit of Great
Britain}: by tying the money supply to the \emph{nominal} (price times quantity)
rather than the \emph{real} volume of trade, the doctrine set up a positive
feedback loop from prices to money creation to prices.  If the central bank
accommodates every demand for credit arising from trade conducted at current prices,
and if prices are rising, the money supply will expand along with prices---adding
fuel to inflation rather than quenching it.

Despite Thornton's critique, the doctrine retained influential adherents.  Thomas
Tooke and John Fullarton rehabilitated it in a weaker form in the mid-nineteenth
century currency--banking debate, emphasizing the ``law of reflux'': notes lent
out against real bills flow back to the issuing bank when the bills mature and are
self-liquidating, so the supply of notes cannot permanently exceed the demand.

By the time the Federal Reserve Act was drafted, the real bills doctrine had
achieved near-canonical status among American banking reformers.  As Federal Reserve
Bank of Richmond economist Robert Hetzel (\citeyear{hetzel2014}) documented, ``The
founders of the Federal Reserve desired to end financial panics \ldots\ [T]hey
understood financial panics as resulting from speculative excess, especially on
Wall Street.  These `real bills' views of financial panic originated in the
19th~century American experience.''

\subsection{Congress's Dual Commitment in 1913}

Parintha Sastry (\citeyear{sastry2018}) has argued that in its original passage the
1913~Act reflected a ``steadfast commitment'' to the real bills doctrine visible in
two interrelated ways: first, by limiting the assets the Fed could purchase, discount,
and accept as collateral to eligible commercial paper; second, by keeping
government-sponsored credit enterprises separate from the Federal Reserve System.
The Fed was to be a lender of last resort to \emph{commercial} banking, not a
financier of speculative or investment activity.

This limitation answered the rural critics of the Aldrich Plan who had feared that
a central bank would become the instrument of Wall Street finance.  By
confining eligible paper to short-term, self-liquidating commercial loans, the
the drafters---Congressman Carter Glass and Senator Robert Latham Owen---sought to
ensure that Fed credit would flow to Main Street merchants and farmers, not to
stock-market speculators.  Wilson secured the passage of the bill in the Senate by
convincing Bryan's supporters that Federal Reserve notes, being obligations of the
government, would satisfy the Populist demand for an ``elastic currency'' without
replicating either the bond-backed rigidity of the National Banking System or the
Wall Street favoritism of the Aldrich Plan.

\subsection{The Echo of Earlier Reform Proposals}

It is worth pausing to hear the echo of the pre-1908 reform proposals in the
structure of the 1913~Act.  The asset-currency advocates of the 1880s and 1890s
had argued that notes should be backed by productive assets rather than government
bonds; by 1913, the drafters went further---they required not merely that bank
assets \emph{back} the currency but that the specific assets eligible for Fed
discount be active claims on ongoing commercial production.  The warehouse-receipt
proposals of agrarian reformers found a close analogue in the six-month agricultural
paper carve-out in Section~13.  The railroad-bond proposals did \emph{not} survive:
investment securities of any kind were explicitly excluded from eligible paper.
The final legislation represented a deliberately narrow conception of what the
currency should be ``backed by''---narrower even than the Aldrich--Vreeland
emergency currency, which had allowed almost any security as backing.

%% ------------------------------------------------------------------
\section{Post-1913 Controversies: Real Bills, Elasticity, and Instability}
\label{sec:aftermath}
%% ------------------------------------------------------------------

\subsection{The Inflationary Bias}

Critics of the real bills doctrine pointed to its inflationary potential almost
from the opening of the Federal Reserve banks.  James Parthemos of the
Federal Reserve Bank of Richmond wrote in 1988:
\begin{quote}
``This so-called commercial loan theory or real bills doctrine was a basic principle
underlying the money functions of the new system.  The essential fallacy in the
doctrine was that note issue would also vary with the price level as well as the
real volume of trade.  Thus its operation would be inherently inflationary or
deflationary.'' \citep{parthemos1988}
\end{quote}
Friedman and Schwartz (\citeyear{friedman-schwartz1963}) noted the same flaw, though
without exploring its full implications.

The pro-cyclical logic was stark.  In a boom, prices rise, the nominal value of
commercial transactions increases, demand for discount rises, and the Fed---if
mechanically following the real bills criterion---expands credit.  In a contraction,
the reverse occurs, amplifying the downturn.

\subsection{The Great Depression and ``Direct Pressure''}

The most consequential invocation of the real bills doctrine came in late 1929.
Long-serving Federal Reserve Board member Adolph C.\ Miller launched what Humphrey
and Timberlake (\citeyear{humphrey-timberlake2019}) called the ``Direct Pressure''
initiative: a directive to all member banks requiring any bank seeking discount
window assistance from the Fed to certify that it had never made speculative
loans---especially loans for stock-market purposes.  Under the real bills doctrine,
speculative finance was categorically different from legitimate commercial lending,
and the Fed was to be the commercial system's lender of last resort, not a
backstop for speculators.

The result was catastrophic.  Bankers, unwilling to submit to the implied
interrogation (and doubting that the ``hard-boiled'' Fed would grant relief
anyway), refrained from borrowing reserves even in the face of withering deposit
outflows.  Instead of discounting at the window, banks failed---nearly 9{,}000 of
them.  Their failure produced the one-third contraction of the money stock that
Friedman and Schwartz held responsible for the Great Depression.

\subsection{The Sargent--Wallace Reconsideration}

In a 1982 article, Thomas Sargent and Neil Wallace (\citeyear{sargent-wallace1982})
reexamined the real bills doctrine in the context of overlapping generations models.
They showed that within such models the real bills prescription---unfettered private
intermediation---could be Pareto superior to the quantity-theoretic prescription,
which imposes restrictions on private intermediation designed to separate ``money''
from credit.  David Laidler (\citeyear{laidler1984}) countered that the Sargent--Wallace
version of real bills differed substantially from the historical doctrine that had
guided Fed policy: the historical real bills criterion was specifically about the
\emph{type} of collateral (short-term commercial paper versus speculative finance),
not simply about whether intermediation should be restricted.

\subsection{Legacy}

The real bills provisions of the original Federal Reserve Act were gradually
eroded over the following decades.  The Banking Acts of 1932 and 1933, passed
during the Depression, expanded the range of assets the Fed could discount; the 1935
amendments transferred power decisively to the Board of Governors and weakened the
operational role of the reserve banks in policing eligible paper.  Section~13(3),
added to the Federal Reserve Act in 1932 and invoked extensively during the
2007--2009 financial crisis, authorizes the Fed to lend to any ``individual,
partnership, or corporation'' in ``unusual and exigent circumstances.''
The provision traces its legislative DNA to the Aldrich--Vreeland emergency
currency mechanism of 1908, and stands at the furthest possible remove from the
strict real bills criterion of the original Act.

%% ------------------------------------------------------------------
\section{Conclusion}
\label{sec:conclusion}
%% ------------------------------------------------------------------

The line from the National Banking Acts of 1863--1864 to the Federal Reserve Act
of 1913 is coherent and traceable.  The bond-backed currency of the national
banking era was safe but rigid; the Comptroller of the Currency guaranteed the safety
by inspecting bond portfolios, but no institutional mechanism existed to make
the currency responsive to the needs of trade.  Four decades of complaint---from
Greenbackers, silver men, asset-currency advocates, and mainstream bankers
alike---converged on the diagnosis of ``inelasticity.''

The Aldrich--Vreeland Act of 1908 was a halfway house: it broadened the
permissible backing for emergency notes beyond government bonds, pointed toward the
eventual solution, and created the intellectual infrastructure (the National Monetary
Commission) that would produce the 1913~Act.  The Federal Reserve Act then
resolved the inelasticity problem by substituting the \emph{real bills} criterion for
the bond-backing requirement.  Instead of asking ``Is this note backed by a government
bond?'' the new regime asked ``Is this note backed by a short-term, self-liquidating
commercial claim?''

The real bills criterion was simultaneously a monetary theory (the money supply
ought to expand and contract with productive activity), a political compromise
(Fed credit should serve Main Street, not Wall Street), and an echo of the long
debate over what money should be ``backed by.''  Its subsequent history---above all
its disastrous role in the onset of the Great Depression---showed that the doctrine
embedded in 1913 carried its own serious defects.  But understanding how the
doctrine came to be embedded in the Act requires tracing the full arc of American
monetary thought from the Civil War to the Progressive Era, a story whose
protagonists include not only the drafters of legislation but the farmers, bankers,
and economists whose decades of argument about ``currency inelasticity'' shaped the
institutional response.

%% ------------------------------------------------------------------
\bibliographystyle{plainnat}
\begin{thebibliography}{99}

\bibitem[Friedman and Schwartz(1963)]{friedman-schwartz1963}
  Friedman, Milton, and Anna~J. Schwartz.
  \newblock \emph{A Monetary History of the United States, 1867--1960}.
  \newblock Princeton University Press, Princeton, N.J., 1963.

\bibitem[Hetzel(2014)]{hetzel2014}
  Hetzel, Robert~L.
  \newblock The real bills views of the founders of the Fed.
  \newblock \emph{Economic Quarterly}, Federal Reserve Bank of Richmond,
    Second Quarter 2014.

\bibitem[Humphrey and Timberlake(2019)]{humphrey-timberlake2019}
  Humphrey, Thomas~M., and Richard~H. Timberlake.
  \newblock \emph{Gold, the Real Bills Doctrine, and the Fed: Sources of
    Monetary Disorder, 1922--1938}.
  \newblock Cato Institute, Washington, D.C., 2019.

\bibitem[Johnson(1999)]{johnson1999}
  Johnson, Roger~T.
  \newblock \emph{Historical Beginnings \ldots\ The Federal Reserve}.
  \newblock Federal Reserve Bank of Boston, Boston, 1999.

\bibitem[Laidler(1984)]{laidler1984}
  Laidler, David.
  \newblock Misconceptions about the real-bills doctrine: A comment on Sargent
    and Wallace.
  \newblock \emph{Journal of Political Economy}, 92\penalty0 (1):\penalty0
    149--155, February 1984.

\bibitem[Laughlin(1908)]{laughlin1908}
  Laughlin, J.~Laurence.
  \newblock The Aldrich-Vreeland Act.
  \newblock \emph{Journal of Political Economy}, 16\penalty0 (8):\penalty0
    489--513, 1908.

\bibitem[Law(1705)]{law1720}
  Law, John.
  \newblock \emph{Money and Trade Considered: With a Proposal for Supplying a
    Nation with Money}.
  \newblock Andrew Anderson, Edinburgh, 1705.

\bibitem[Mints(1945)]{mints1945}
  Mints, Lloyd~W.
  \newblock \emph{A History of Banking Theory in Great Britain and the United
    States}.
  \newblock University of Chicago Press, Chicago, 1945.

\bibitem[Parthemos(1988)]{parthemos1988}
  Parthemos, James.
  \newblock The Federal Reserve Act of 1913 in the stream of U.S.\ economic
    history.
  \newblock \emph{Economic Review}, Federal Reserve Bank of Richmond,
    July/August 1988.

\bibitem[Sargent and Wallace(1982)]{sargent-wallace1982}
  Sargent, Thomas~J., and Neil Wallace.
  \newblock The real-bills doctrine versus the quantity theory: A
    reconsideration.
  \newblock \emph{Journal of Political Economy}, 90\penalty0 (6):\penalty0
    1212--1236, December 1982.

\bibitem[Sastry(2018)]{sastry2018}
  Sastry, Parintha.
  \newblock The political origins of Section~13(3) of the Federal Reserve Act.
  \newblock \emph{Economic Policy Review}, Federal Reserve Bank of New York,
    2018.

\bibitem[Silber(2007)]{silber2007}
  Silber, William~L.
  \newblock The great financial crisis of 1914: What can we learn from
    Aldrich-Vreeland emergency currency?
  \newblock \emph{American Economic Review (Papers and Proceedings)},
    97\penalty0 (2):\penalty0 285--289, 2007.

\bibitem[Smith(1776)]{smith1776}
  Smith, Adam.
  \newblock \emph{An Inquiry into the Nature and Causes of the Wealth of
    Nations}.
  \newblock Strahan and Cadell, London, 1776.

\bibitem[Thornton(1802)]{thornton1802}
  Thornton, Henry.
  \newblock \emph{An Enquiry into the Nature and Effects of the Paper Credit of
    Great Britain}.
  \newblock Hatchard, London, 1802.

\bibitem[Wicker(2005)]{wicker2005}
  Wicker, Elmus.
  \newblock \emph{The Great Debate on Banking Reform: Nelson Aldrich and the
    Origins of the Fed}.
  \newblock Ohio State University Press, Columbus, 2005.

\bibitem[Willis(1914)]{willis1914}
  Willis, Henry~Parker.
  \newblock The Federal Reserve Act.
  \newblock \emph{American Economic Review}, 4\penalty0 (1):\penalty0 1--24,
    1914.

\end{thebibliography}

\end{document}
